% SPDX-License-Identifier: CC-BY-4.0
\section{The Hyphae Architecture}
\label{sec:arch}

\hyphae{} is structured as three concentric layers:
\textbf{substrate} (memory state, fragment storage,
hash-chained journal), \textbf{cognition path} (the operations
that retrieve, compose, and emit), and \textbf{realizer}
(the prose-emission stage). The architectural commitments
load-bearing for the empirical results are concentrated at the
realizer boundary; we describe each layer briefly with that
focus.

\subsection{Substrate}

The substrate stores \emph{cognitive fragments}: small
self-contained units of content with provenance metadata
(source subsystem, parent fragments if cascade-derived,
confabulation risk score, valence, domain tags). Fragments
are written to disk via an embedded key-value store
\citep{redb, fjall} and chained into a SHA-256 hash sequence
that links each fragment to its predecessor by content hash.
The journal is verifiable: any modification to a previously-
written fragment breaks the chain at and after the modification
point.

Retrieval is via approximate-nearest-neighbour search
\citep{malkov2018hnsw} over fragment embeddings produced by a
deterministic hashing token embedder. The embedder is not an
LLM, contains no learned parameters, and uses the
\texttt{HashingTokenEmbedder} primitive from the substrate.
This makes the retrieval stage deterministic and reproducible
across runs.

\subsection{Cognition Path}

A \textbf{cognition path} is a procedure that reads from the
substrate, performs a structured operation, and either writes
new fragments back or emits a response through the realizer.
The v0.1 substrate implements five paths: \texttt{ingest}
(write a new fragment), \texttt{recall} (retrieve fragments
similar to a cue), \texttt{compose} (synthesize a response from
retrieved fragments), \texttt{learning\_update} (adjust
substrate-internal parameters under audit), and
\texttt{grounded\_retrieval} (deferred per the v0.1 RFC).

The \texttt{compose} path is the one this paper measures. It
takes a \emph{working set} of cognitive fragments (typically
retrieved by a \texttt{recall} call) and produces a
\texttt{RealizationOutput}: composed prose, a list of
fragment IDs quoted in emission order, and a structured set of
limitation triggers (such as \texttt{HighConfabRisk} or
\texttt{ShallowCascade}) that fired during composition.

\subsection{Realizer}
\label{sec:realizer}

The realizer is the surface-prose-emission stage. It is the
component where the LLM would conventionally sit. In \hyphae{}
it is a finite procedure:

\begin{enumerate}
    \item \textbf{Schema selection.} The caller's intent (one
    of: dialogue, assert, summarize, compare, reflect, narrate)
    maps deterministically to a schema (DialogueReply,
    GroundedAssertion, Summary, ComparativeAnalysis,
    IntrospectiveAssessment, NarrativeArc).
    \item \textbf{Shape derivation.} If the caller supplies an
    explicit composition shape, the realizer walks its steps.
    Otherwise the realizer derives a shape from the working set
    by detecting opposed valence (Contrast role) and adjacency
    relations (Continuation role). The shape determines the
    sequence of connective roles between adjacent quotations.
    \item \textbf{Connective selection.} For each inter-fragment
    boundary, the realizer queries the lexicon for a connective
    matching the (role, register, polarity, formality) tuple.
    Boundary smoothing rules (Section~\ref{sec:smoothing})
    filter out connectives whose head would clash with the
    adjacent quote's tail.
    \item \textbf{Quote emission.} Each fragment's body is
    emitted verbatim inside double quotes. No paraphrase, no
    summarisation, no reordering of intra-body text.
    \item \textbf{Limitation acknowledgement.} Triggers that
    fired during composition (empty working set, high
    confabulation risk, shallow cascade depth, ethically
    sensitive content) are surfaced via the lexicon's
    limitation-acknowledgement phrases.
\end{enumerate}

The realizer's output is bounded: every emitted token comes
either from the curated lexicon (approximately 250 phrases
covering 10 connective roles, 4 conversational registers, 4
polarity classes, and 3 formality tiers) or from a fragment
body quoted verbatim. No statistical model decides the
emission.

\subsection{Hard Commitment 12 --- Verbatim Quotation}
\label{sec:hc12}

The architectural decision underlying the empirical results
is \emph{Hard Architectural Commitment 12} from the project's
foundational ADR (\texttt{docs/adr/0001-fresh-from-v1.md}):

\begin{quote}
\emph{Composition uses fragment quotation plus connective tissue,
not novel language synthesis. Fragments are opaque content
sources whose body text is preserved verbatim. The surface
realizer generates only the structural prose connecting them.
This boundary is load-bearing for the no-LLM-in-cognition-path
claim.}
\end{quote}

This commitment trades expressive flexibility for two
properties that anchor the empirical comparison in
Section~\ref{sec:results}:

\textbf{Verifiability.} Every quoted body in a \hyphae{} response
is byte-identical to a fragment stored in the substrate, which
is itself verifiable against the hash-chained journal. The
audit relation is total: a third party can independently confirm
that the output was constructed from named substrate fragments
without re-running the system.

\textbf{Unsupported-claim immunity by construction.} The NLI
metric we employ (Section~\ref{sec:metrics}) cannot mark a
verbatim quote as unsupported with respect to its source: the
quote and the context are identical text. The connective tissue
between quotes is non-factual scaffolding the metric's filter
recognises and excludes from the denominator. The result is that
on factual-retrieval tasks, the architectural commitment maps
directly onto the metric's scoring shape.

\subsection{Lexicon and Boundary Smoothing}
\label{sec:smoothing}

The lexicon is a curated mapping from
\emph{(role, register, polarity, formality)} tuples to surface
phrases. The baseline EN lexicon contains approximately 250
entries organised across 10 connective roles
(Opening, Continuation, Contrast, Attribution, Closing,
Concession, Causation, Elaboration, Sequence, Summary), 4
register classes (Neutral, Formal, Conversational, Technical),
4 polarity classes (Continuation, ContrastSoft, ContrastHard,
Concession, Neutral), and 3 formality tiers (Low, Mid, High).

The picker has a four-level fallback chain: an exact
$(role, register, polarity, formality)$ tuple match is preferred;
on miss, formality relaxes, then register, then polarity, until
finally any phrase for the requested role is returned.

\textbf{Boundary smoothing} avoids surface-level disfluencies at
the boundary between a connective and an adjacent fragment quote.
For example, the connective \emph{``Per the recorded fragments,''}
followed by a quote that begins with \emph{``Per the''} would
produce a stuttered repetition. The smoothing pass extracts
boundary signals from the adjacent quote (leading determiner,
anaphor surface form) and filters the candidate connective set
to avoid such collisions.

Section~\ref{sec:ablations} reports ablations of cascade-shape
composition, ethics gate, lexicon scale, and boundary smoothing.
Of these, only lexicon scale produces a measurable effect on the
comparator metrics --- a finding that constrains claims about
the contribution of the other three components at our corpus
size.

\subsection{What the Realizer Does Not Do}

The realizer does not paraphrase, summarise, reorder
intra-body text, infer over multiple fragments, or generate
prose outside the lexicon's curated set. These restrictions
are not implementation gaps; they are architectural decisions.
The choice trades prose-fluency dimensions in which an LLM
excels for the dimensions catalogued above.

For applications requiring paraphrastic flexibility or
synthesis over multiple sources, \hyphae{}'s realizer is not
the right tool. The right comparator question is therefore
not ``does \hyphae{} replace an LLM,'' but rather ``for the
subset of tasks where verbatim grounding suffices, is the
substrate competitive?''. Sections~\ref{sec:results}
and~\ref{sec:discussion} take up that question.
